% The component: what it does, its units, interfaces and resources.
% Identification data comes from metadata.tex.
% Example: components/sample/component.tex.

% \FunctionalDescription{
%   \TheComponent\ ...
% }

% Units of the component: {unit}{responsibility}{unit verification}
% \Unit{}{}{}

% Requirements realised by each unit: {unit}{REQ, REQ}
% \Allocate{}{}

% Provided interfaces: [unit that realises it]{name}{type / range / timing}{description}
% \ProvidedInterface[]{}{}{}

% Required interfaces: [providing component]{name}{description}{replaced by}{behaviour in test}
% [component] is optional: interfaces of one component share a box in the diagrams.
% \RequiredInterface[]{}{}{Stub}{}

% Use relations for the structure diagram: {unit}{unit or component}{interfaces}
% \Uses{}{}{}

% Who uses the component: {component}{interfaces it calls}
% \Caller{}{}

% Resource budgets: {resource}{budget}{verified by}
% \Resource{RAM}{}{}

% Configuration parameters: {name}{values / default}{description}
% \ConfigParameter{}{}{}

% Free text: any LaTeX, printed at the end of this chapter. Source files of the
% component go to code/ and are listed with \InputCode{file}.
% \begin{FreeText}
% \subsection{Public API (excerpt)}
% \InputCode{component.h}
% \end{FreeText}
